\qznutitle{泉州世遗文化融入初中语文综合性学习研究}
\qznusubtitle{——基于项目化学习的行动研究}
\qznutitleen{Integrating Quanzhou World Heritage Culture into Junior High School Chinese Comprehensive Learning: An Action Research Based on Project-Based Learning}
\qznucollege{文学与传播学院}
\qznucollegeen{School of Literature and Communication}
\qznumajor{汉语言文学}
\qznumajoren{Chinese Language and Literature}
\qznuclass{2022级1班}
\qznustudentid{2022010101}
\qznuauthor{陈晓岚}
\qznuauthoren{CHEN Xiaolan}
\qznuadvisor{林文远}
\qznuadvisoren{LIN Wenyuan}
\qznuadvisortitle{副教授}
\qznuadvisortitleen{Associate Professor}
\qznudate{2026年5月}
\qznuabstracttext{在课程思政与中华优秀传统文化教育协同推进的背景下，地方世界遗产如何转化为语文课程资源，是初中语文综合性学习亟需回应的问题。本研究以泉州世界遗产为文化资源，以项目化学习为组织方式，在泉州市某中学八年级开展两轮行动研究。通过课堂观察、学生作品分析、问卷调查与半结构化访谈，考察该模式对学生在文化理解、语言运用、合作探究与问题解决等方面的影响。研究发现，项目化学习能够提升学生参与综合性学习的主动性，促进学生在真实情境中整合阅读、表达与探究能力；同时，教师需要在任务支架、文化材料筛选和评价量表设计上提供更精细的支持。}
\qznukeywords{泉州世界遗产；初中语文；综合性学习；项目化学习}
\qznuabstractentext{Against the background of curriculum-based ideological and political education and Chinese traditional culture education, how to transform local world heritage into Chinese curriculum resources is an urgent issue for comprehensive learning in junior high school. Taking Quanzhou World Heritage as cultural resources and project-based learning as the organizational approach, this study conducted two rounds of action research in Grade 8 of a middle school in Quanzhou. Through classroom observation, analysis of student works, questionnaires and semi-structured interviews, it examined the effects of the model on students' cultural understanding, language use, collaborative inquiry and problem solving. The findings show that project-based learning can enhance students' initiative in comprehensive learning and promote their ability to integrate reading, expression and inquiry in authentic contexts. Meanwhile, teachers need to provide more refined support in task scaffolding, selection of cultural materials and design of assessment rubrics.}
\qznukeywordsen{Quanzhou World Heritage; junior high school Chinese; comprehensive learning; project-based learning}
